\chapter{Zusammenfassung und Auswertung}

Die vorliegende Arbeit untersucht die Frage, ob und mit welcher Güte sich Log-Nachrichten kontinuierlich und ohne manuellen Pflegeaufwand gruppieren lassen. Ausgangspunkt ist die in Kapitel~\ref{chap:einf} dargelegte Beobachtung, dass Log-Nachrichten moderner Software-Systeme keinem einheitlichen Format folgen und dass klassische Ansätze wie regelbasierte Regex-Filter einen hohen manuellen Pflegeaufwand erfordern. Gesucht wurden daher Verfahren, die eingehende Log-Nachrichten automatisch Gruppen -- Clustern -- zuordnen und diese Zuordnung bei neuen Nachrichten fortlaufend aktualisieren, ohne manuellen Eingriff zu erfordern.

Zur Umsetzung dieses Ziels wurden zwei grundlegende Verfahrensklassen betrachtet. Kapitel~\ref{chap:ml} behandelt algorithmische Ansätze des maschinellen Lernens: Der partitionierende Algorithmus K-Means sowie der dichtebasierte Algorithmus DBSCAN operieren auf TF-IDF-Vektoren (Term Frequency-Inverse Document Frequency) und erfordern eine vollständige Neuverarbeitung bei jedem Durchlauf. Ergänzend werden mit Drain3 und Brain zwei spezialisierte Online-Parser betrachtet, die Log-Nachrichten iterativ einem wachsenden Template-Baum zuordnen und somit nativ für den kontinuierlichen Betrieb ausgelegt sind. Kapitel~\ref{chap:vek} stellt den vektorbasierten Ansatz vor: Jede Log-Nachricht wird mittels des vortrainierten Sentence-Transformer-Modells \texttt{all-MiniLM-L6-v2} in einen 384-dimensionalen semantischen Vektor überführt und in einer PostgreSQL-Datenbank mit \texttt{pgvector}-Erweiterung persistiert. Ein zweistufiger Clustering-Algorithmus ordnet neu eingehende Nachrichten zunächst bestehenden Clustern zu (Re-Clustering) und legt nur bei semantischer Neuartigkeit einen neuen Cluster an. Als weiterer Vergleichskandidat wird MinHash LSH (Locality-Sensitive Hashing) untersucht, das auf Signaturen statt auf Vektoren operiert.

Die praktische Realisierung aller Verfahren erfolgte in einer reproduzierbaren Docker-Umgebung bestehend aus einem Jupyter-Notebook-Container und einer PostgreSQL-Instanz mit \texttt{pgvector}-Erweiterung, wie in Kapitel~\ref{chap:real} beschrieben. Als Datenbasis dienen die drei Systemlog-Quellen \texttt{Linux}, \texttt{Apache} und \texttt{OpenSSH} aus dem Loghub-2k-Benchmark \cite{zhu2023loghublargecollectionlog}, die zu einem gemeinsamen Datensatz mit $10\,000$ annotierten Log-Nachrichten und $44$ Ground-Truth-Klassen zusammengeführt werden. Jeder Algorithmus wurde in fünf Durchläufen mit zufällig permutierter Eingangsreihenfolge evaluiert, um das Verhalten im kontinuierlichen Betrieb abzubilden. Die Clustering-Güte wird anhand von zwei komplementären Metriken beurteilt: dem Adjusted Rand Index~(ARI), der die Übereinstimmung mit der Ground Truth zufallskorrigiert misst, und dem V-Measure, das Homogenität und Vollständigkeit der gefundenen Cluster als harmonisches Mittel vereint.

Das folgende Kapitel fasst die Ergebnisse dieser Versuche vergleichend zusammen, diskutiert die Stärken und Schwächen der einzelnen Verfahren und leitet daraus Empfehlungen für den produktiven Einsatz ab.

\section{Vergleichende Auswertung der Clustering-Verfahren}

\section{Weitere Ansätze}



\section{Nächste Schritte}

\section{Ausblick}
